\subsection{Harmonic residue characters and the exact remaining condition}

Here \(B_g\) denotes the full adelic Alexander unit evaluated by a
finite character at each coefficient prime. Put \(G=\widehat{\mathrm{GT}}\),
\(H=G_{\mathbf Q}\), \(N=\ker(\lambda,\rho_2)\), and
\(D=N_{\rm rad}\). For an odd prime \(\ell\), choose \(i\) in
\(\overline{\mathbf F}_\ell\) with \(i^2=-1\), and put
\[
 h_\ell=\frac{\ell-\chi_{-4}(\ell)}2,\qquad
 (-1)^{\vartheta_{\ell,4}(g)}
       =B_{g,\ell}(i,-1)^{h_\ell}\quad(g\in N).
\]
The value is a sign, independent of the choice of \(i\).
Theorem~\ref{thm:all-kummer-parity} gives the evenness of every
\(\rho_d^I\) on multiplier one, so the divided Kummer coordinate
in the following statement is defined on all of \(N\).

\begin{proposition}\label{prop:harmonic-finite-relative-character}
The map \(\vartheta_{\ell,4}:N\to\mathbf F_2\) is a continuous,
\(G\)-invariant homomorphism. The difference
\[
 \Xi_{\ell,4}(g)=\vartheta_{\ell,4}(g)
                    +\frac{\rho_\ell^I(g)}2\bmod2
                   \qquad(g\in N)
\]
vanishes on \(N\cap H\), and therefore extends uniquely
to a continuous character \(G\to\mathbf F_2\) that is zero
on \(H\).
On \(D\) this character is \(\vartheta_{\ell,4}\).
By Theorem~\ref{thm:harmonic-beta-square-lattice}, its vanishing
there is exactly the tame harmonic square test at \(\ell\).
No nonvanishing of \(\Xi_{\ell,4}\) on full GT is asserted.
\end{proposition}
\begin{proof}
On the multiplier-one group, the beta cocycle is
multiplicative. The unitary identity gives
\[
 B_{g,\ell}(i,-1)B_{g,\ell}(-i,-1)=1.
\]
At an inert prime, Frobenius exchanges these values, so the
first lies in the norm-one group of order \(\ell+1\).
At a split prime the pair identifies the split norm-one
torus with \(\mathbf F_\ell^\times\), of order \(\ell-1\).
Its \(h_\ell\)-th power is consequently a sign and is
unchanged by inversion. This proves the homomorphism and
orientation assertions.
The semilinear beta law gives, for \(a\in G\) and \(g\in N\),
\[
 B_{aga^{-1}}(X,Y)=B_g(X^{\lambda_a},Y^{\lambda_a}).
\]
At \((i,-1)\), the only possible change is \(i\mapsto-i\);
the sign is unchanged. Thus \(\vartheta_{\ell,4}\) is invariant.
The divided Kummer character is invariant for the same
reason: conjugation scales \(\rho_\ell^I|_N\) by the odd
unit \(\lambda_a\).

We check the arithmetic restriction with the full Anderson
formula, not only its logarithmic derivative. In the beta
form of \cite[Remark~2.4, p.~101]{Furusho2003}, put
\(n=\ell\) and \(s=1/4\), and reduce in characteristic
\(\ell\). Every \(\ell\)-th root in the denominator reduces
to one; the denominator is therefore one.
Normalized duplication gives
\[
 B_\sigma(i,i)=B_\sigma(-1,i)=B_\sigma(i,-1)
                     \quad(\sigma\in N\cap H).
\]
In the numerator, the residues \(k=1,2\bmod4\) contribute
this common value and the residues \(k=0,3\bmod4\)
contribute one. Their total number for
\(0\le k\le\ell-1\) is exactly \(h_\ell\):
it is \(2r\) if \(\ell=4r+1\), and \(2r+2\) if
\(\ell=4r+3\). The arithmetic right side is
\[
 i^{-\ell\rho_\ell^I(\sigma)}
       =(-1)^{\rho_\ell^I(\sigma)/2},
\]
since the Kummer coordinate is even on the cyclotomic
kernel. This proves \(\Xi_{\ell,4}|_{N\cap H}=0\).

Finally \(G=NH\). The invariant character extends by
\(\Xi_{\ell,4}(nh)=\Xi_{\ell,4}(n)\). Vanishing on the intersection
proves that this is well-defined, invariance proves
multiplicativity, and compactness of \(N\times H\to G\)
proves continuity. Its uniqueness follows from \(G=NH\).
\end{proof}

There is a finite-cover interpretation of the beta
coordinate used here. The cyclic degree-four cover
\[
                 v^4=t(t-1)^2
\]
has monodromy \(x\mapsto1,\ y\mapsto2\) in
\(\mathbf Z/4\). Its compactification is the elliptic curve
\[
             v^2=z^3-z,\qquad t=z^2,\quad
             z=\frac{v^2}{t-1}.
\]
The \((i,-1)\) eigenspace in its homology modulo \(\ell\)
is one-dimensional, and the restricted GT action
on that eigenspace is \(B_{g,\ell}(i,-1)\).
Indeed the three nontrivial deck characters occur once in
the homology of the punctured cover; the order-two character
is a boundary character, and the two primitive characters
form the closed elliptic homology. The Alexander generator
maps nontrivially to each primitive eigenspace, so its
coefficient is exactly the stated beta value.
This interpretation uses restriction of an automorphism
that preserves the displayed subgroup. It makes no
assertion that every possible scalar lifts to a GT element.

\begin{proposition}\label{prop:harmonic-formal-residue-gap}
Vanishing of the integral rational Kummer coordinates,
beta associativity, unitarity, the dyadic Anderson formulas,
and logarithmic Anderson identities do not by themselves
force every harmonic square test as identities of formal
integral Alexander data.
\end{proposition}
\begin{proof}
Fix an odd prime \(\ell\ge5\). The established simultaneous
arithmetic radical image supplies
\(\sigma\in H^{\rm cyc}\) with
\[
 \rho_\ell^I(\sigma)=2e_2,\qquad
 \rho_q^I(\sigma)=0\quad(q\ne\ell),
\]
where \(e_2\in\widehat{\mathbf Z}\) is the idempotent
supported at the coefficient prime two. In particular
\(\sigma\in N\) and every coefficient-\(\ell\) Kummer
coordinate is zero. The preceding arithmetic calculation
gives
\[
              B_{\sigma,\ell}(i,-1)^{h_\ell}=-1.
\]
Keep only the coefficient-\(\ell\) component of
\(R_\sigma\) and set its other coefficient-prime components
to zero. The resulting beta unit is \(B_{\sigma,\ell}\)
in this component and one in all others.
All integral cyclic-cover coefficient maps are
\(\widehat{\mathbf Z}\)-linear on the Alexander module.
Thus every patched Kummer coordinate is zero.
Associativity, unitarity and normalized dyadic identities
hold componentwise. For an Anderson discrepancy near a
prime-to-\(\ell\) torsion point, its arithmetic factor is
a constant root of unity times the principal factor
\((1+S)^{n\rho_{n,\ell}^I(\sigma)}\).
The latter is one, so its logarithmic derivative vanishes.
At the other coefficient primes the beta unit is one.
The patched data therefore satisfy all the listed formal
conditions while retaining the nontrivial harmonic sign.

This is formal Alexander data. No lift satisfying the full
profinite pentagon, or the full nonabelian cyclic-filling
identities, is constructed. The construction isolates a
missing implication and is not a counterexample in GT.
\end{proof}
